% =============================================================================
% Chapter 2 — The Stack ($0100–$01FF)
% =============================================================================
\chapter{The Stack}

\epigraph{``Push it down, pull it back up.  If only life were so
orderly.''}{\textit{--- Anonymous assembly programmer, ca.\ 1979}}

\section*{Introduction}

On the MOS 6502, the hardware stack occupies exactly one page of memory: the
256 bytes from \$0100 to \$01FF.  There is no option to move it, extend it, or
protect it.  The processor's 8-bit stack pointer (SP) is simply an offset into
page 1 --- push an item and SP decrements; pull an item and SP increments.
When SP wraps from \$00 back to \$FF (or vice versa), the CPU does not complain.
It just keeps going, silently overwriting whatever was there before.

This simplicity is both a strength and a hazard.  The stack is fast and
predictable, but 256 bytes is not very much, especially when BASIC is using it
for subroutine calls, loop variables, and expression evaluation all at once.

\memrange{0100}{01FF}{Hardware Stack}{R/W}{---}

The entire 256-byte region is ordinary read/write RAM.  The CPU treats it as
special only because push and pull instructions always address page 1 using SP
as the low byte.  Nothing prevents you from reading or writing stack memory
directly with PEEK and POKE --- just be aware that the interpreter is using it
too.

% ═══════════════════════════════════════════════════════════════════════════════
\section{Stack Mechanics}
% ═══════════════════════════════════════════════════════════════════════════════

The 6502 stack grows \emph{downward} in memory.  When the system starts, SP is
initialised to \$FF, meaning the first byte pushed will be written to \$01FF.
After the push, SP decrements to \$FE.  The next push writes to \$01FE, and so
on.

Pulling works in reverse: SP increments first, then the byte at the new
address is read.  This means the ``current'' top of stack is always at
\texttt{\$0100 + SP + 1} --- the byte just above where SP points.

\begin{center}
\small
\begin{tabular}{@{}lll@{}}
\toprule
\textbf{Operation} & \textbf{Effect on SP} & \textbf{Address Used} \\
\midrule
\texttt{PHA} / \texttt{PHP}     & SP $\leftarrow$ SP $-$ 1 & Write to \$0100 + SP (before decrement) \\
\texttt{PLA} / \texttt{PLP}     & SP $\leftarrow$ SP $+$ 1 & Read from \$0100 + SP (after increment) \\
\texttt{JSR}                     & SP $\leftarrow$ SP $-$ 2 & Pushes PCH then PCL \\
\texttt{RTS}                     & SP $\leftarrow$ SP $+$ 2 & Pulls PCL then PCH \\
\texttt{BRK} / IRQ / NMI        & SP $\leftarrow$ SP $-$ 3 & Pushes PCH, PCL, then P \\
\texttt{RTI}                     & SP $\leftarrow$ SP $+$ 3 & Pulls P, then PCL, then PCH \\
\bottomrule
\end{tabular}
\end{center}

\begin{warningbox}
The 6502 performs no overflow or underflow detection on the stack pointer.
If SP wraps from \$00 to \$FF during a push, the CPU will happily write into
the top of the stack page, overwriting data that was pushed earlier.
Similarly, pulling past \$FF wraps to \$00.  There is no trap, no error
message --- just silent data corruption.
\end{warningbox}

% ═══════════════════════════════════════════════════════════════════════════════
\section{Stack Frame Formats}
% ═══════════════════════════════════════════════════════════════════════════════

The stack is used by both the 6502 hardware and EhBASIC for several distinct
frame types.  Understanding these formats helps when debugging programs that
crash or behave unexpectedly.

\subsection{JSR/RTS --- Subroutine Call (2 bytes)}

When the CPU executes a \texttt{JSR addr} instruction, it pushes the return
address \emph{minus one} onto the stack, high byte first, then low byte:

\begin{center}
\begin{tabular}{@{}ll@{}}
\toprule
\textbf{Stack offset} & \textbf{Content} \\
\midrule
SP + 2 & PCH (high byte of return address $-$ 1) \\
SP + 1 & PCL (low byte of return address $-$ 1) \\
\bottomrule
\end{tabular}
\end{center}

The ``minus one'' quirk is because \texttt{RTS} increments the pulled address
before jumping to it.  This is a well-known 6502 subtlety and the reason that
\texttt{JMP} and \texttt{JSR} are not interchangeable in dispatch tables.

\subsection{BRK/IRQ/NMI --- Interrupt Frame (3 bytes)}

Hardware interrupts and the \texttt{BRK} instruction push three bytes: the full
return address (high byte first, then low byte), followed by the processor
status register (P):

\begin{center}
\begin{tabular}{@{}ll@{}}
\toprule
\textbf{Stack offset} & \textbf{Content} \\
\midrule
SP + 3 & PCH \\
SP + 2 & PCL \\
SP + 1 & P (processor status) \\
\bottomrule
\end{tabular}
\end{center}

\texttt{RTI} reverses the process: it pulls P, then PCL, then PCH, and resumes
execution at the restored address.  Unlike \texttt{RTS}, \texttt{RTI} does
\emph{not} add one to the return address --- the pushed PC already points to
the correct instruction.

\subsection{EhBASIC FOR/NEXT Frame (\raisebox{0pt}{\texttildelow}18 bytes)}

Each active \texttt{FOR} loop pushes a substantial frame onto the stack so that
\texttt{NEXT} can find the loop parameters:

\begin{center}
\small
\begin{tabular}{@{}lll@{}}
\toprule
\textbf{Offset} & \textbf{Size} & \textbf{Content} \\
\midrule
+0  & 1 byte  & FOR token (\$81) --- identifies this as a FOR frame \\
+1  & 4 bytes & Loop limit (floating-point) \\
+5  & 4 bytes & Step value (floating-point) \\
+9  & 1 byte  & Step sign \\
+10 & 2 bytes & Current line number \\
+12 & 2 bytes & Text pointer (resume position in program text) \\
+14 & 2 bytes & FOR variable pointer (address of loop counter) \\
\bottomrule
\end{tabular}
\end{center}

That is approximately 18 bytes per \texttt{FOR} loop.  The stack can hold at
most about seven nested \texttt{FOR}/\texttt{NEXT} loops before running into
trouble.

\subsection{EhBASIC GOSUB Frame (5 bytes)}

A \texttt{GOSUB} call pushes a much smaller frame:

\begin{center}
\begin{tabular}{@{}lll@{}}
\toprule
\textbf{Offset} & \textbf{Size} & \textbf{Content} \\
\midrule
+0 & 1 byte  & GOSUB token (\$89) \\
+1 & 2 bytes & Current line number \\
+3 & 2 bytes & Text pointer (resume position after RETURN) \\
\bottomrule
\end{tabular}
\end{center}

Five bytes per \texttt{GOSUB} means you can nest about 12 subroutine calls
before the stack fills up.

% ═══════════════════════════════════════════════════════════════════════════════
\section{Stack Pointer Initialisation}
% ═══════════════════════════════════════════════════════════════════════════════

When the NovaVM resets, the CPU fetches the reset vector from \$FFFC--\$FFFD and
jumps to the monitor code at \$FF80.  One of the first things the monitor does
is initialise the stack pointer:

\begin{lstlisting}
LDX #$FF
TXS
\end{lstlisting}

This sets SP to \$FF, placing the stack at the very top of page 1.  EhBASIC
then takes over and begins using the stack for its own purposes.  By the time
you see the \texttt{Ready} prompt, SP is typically somewhere around \$F8 ---
the interpreter keeps a few bytes on the stack at all times for its main loop.

% ═══════════════════════════════════════════════════════════════════════════════
\section{Reading the Stack Pointer}
% ═══════════════════════════════════════════════════════════════════════════════

The 6502's stack pointer is an internal CPU register and cannot be read
directly with PEEK.  However, the \texttt{TSX} instruction copies SP into X,
and a small machine-language routine can make this value available to BASIC:

\begin{lstlisting}
10 REM read current stack pointer
20 REM SP is not directly readable from BASIC,
30 REM but you can watch stack depth via USR()
40 REM or CALL to a small assembly routine
\end{lstlisting}

You can also get a rough sense of stack usage by examining memory near \$01FF
and looking for the boundary between pushed data and uninitialised space.

% ═══════════════════════════════════════════════════════════════════════════════
\section*{Practical Limits}
% ═══════════════════════════════════════════════════════════════════════════════

\begin{tipbox}
In practice, the stack has room for about 12 nested GOSUB calls or 7 nested
FOR/NEXT loops before overflow.  A mix of both will reduce these limits
further.  If your program uses deep recursion via GOSUB, consider restructuring
it with GOTOs or moving the inner logic into machine language.
\end{tipbox}

\begin{warningbox}
Stack overflow in EhBASIC manifests as bizarre behaviour: variables change
unexpectedly, the program jumps to wrong lines, or the interpreter locks up
entirely.  Unlike modern systems, there is no ``stack overflow'' error message.
If your program behaves erratically after adding nested loops or subroutines,
check your nesting depth first.
\end{warningbox}
